\chapter{Conclusiones y perspectivas}
\small
\begin{enumerate}

\item Los parámetros cinemáticos y la recurrencia de CMEs homólogas en la región de origen determinan la viabilidad de una interacción en el medio interplanetario, durante la cual la superposición de sus perfiles de densidad y velocidad puede intensificar considerablemente la geoefectividad del evento resultante.

\item La viabilidad de la interacción entre CMEs homólogas está sujeta a umbrales de velocidad relativa y aceleración inicial que dependen del caso particular, en el Caso 1 la interacción es posible sin que la CME sucesora supere a la precursora ($V_{\text{CME-2}}>0{,}86\,V_{\text{CME-1}}$ con $a_0^{\text{CME-2}}>3{,}23\,a_0^{\text{CME-1}}$), mientras que en el Caso 2 esta condición sí resulta necesaria ($V_{\text{CME-3}}>V_{\text{CME-1}}$ con $a_0^{\text{CME-3}}>0{,}26\,a_0^{\text{CME-1}}$).

\item La sensibilidad de la interacción a los parámetros de desaceleración de la CME sucesora ($a_d$, $\tau_d$) depende de la duración de su fase de aceleración ($\tau_r$), si es corta, el tiempo y la altura de interacción se correlacionan fuertemente con $a_d$ y $\tau_d$ ($R^2 \sim 0{,}89$–$1{,}0$); si es larga, esa correlación se pierde casi por completo ($R^2 \sim 0{,}017$–$0{,}089$) y la geometría del encuentro queda gobernada principalmente por la cinemática de la precursora y el tiempo de retraso entre erupciones.

\item La simulación reproduce cualitativamente rasgos observados en interacciones reales de CMEs, como un frente de baja densidad seguido de una región de alta densidad, una estela de densidad reducida tras el paso de las estructuras, y una escala temporal de contacto entre frentes ($10$–$18$ h, $<30\,R_\odot$), pese a que el modelo no resuelve procesos magnetohidrodinámicos como choques o reconexión.

\end{enumerate}